\chapter{Literature Review and Research Gap}\label{ch:lit}

This chapter reviews the work this project builds on and states precisely where it stops. The
emphasis is on what each line of work establishes and what it leaves unresolved, because the
gap this project addresses is defined by the second of those more than the first.

All twenty-two references cited here are the ones carried by the accompanying manuscript. No
reference has been added for this report without its bibliographic details being checked
against the publisher, proceedings or preprint record.

\section{Real-time degradation from on-chip network activity}

The closest prior work, and the direct motivation for this project, is the study by Behnke et
al.\ \cite{behnke2020}. They instrumented a periodic critical task on an ESP32 and applied IP
packet load while measuring the task's lateness. The degradation they report is modest in
percentage terms but unambiguous: between 2.15\% and 3.8\% of cycles lost for traffic addressed
to a listened port, and between 3.17\% and 6.67\% for an unused port, with latency rising as
the microcontroller approached saturation. Their conclusion is that command-and-control over IP
is unsuitable for hard real-time work on this class of device.

Two aspects of that result deserve attention. The first is that interference is higher for the
\emph{unused} port than for the listened one, which is counter-intuitive until one notices that
it is consistent with cost being incurred below the application: the device pays for frame
reception and demultiplexing whether or not anything is waiting for the data. The second is
that the authors are explicit about the limits of their own mitigations, and it is that honesty
that makes the gap visible.

The same group proposed two mitigations in follow-up work. The first is a priority-aware
multiqueue network interface with per-flow interrupt moderation \cite{behnke2022}; the second is
driver-level early demultiplexing \cite{behnke2023}. Neither is available to a developer
shipping a stock ESP32 today. The multiqueue design was evaluated against a \emph{simulated}
network interface and was never built. The demultiplexing work was prototyped on a Xilinx
Zynq-7000 — a platform an order of magnitude larger than an ESP32 — with measurements taken on
a single core.

The detail that matters most for this project appears in \cite{behnke2023}: the ESP32's second
core is used there to separate the observed processes from the test harness, not as a
mitigation. \textbf{Core pinning is therefore nowhere evaluated as a defence} in this line of
work, despite being the mitigation that developers of these devices are actually told to use.

Fischer et al.\ \cite{fischer2020} examine industrial device architectures under
denial-of-service conditions and reach a compatible conclusion at a coarser granularity: device
architecture, rather than protocol choice, determines whether real-time behaviour survives
network stress. Their analysis is architectural rather than experimental at the task level, so
it frames the problem without constraining the answer for any specific chip.

\section{Timing isolation on larger platforms}

A substantial body of work addresses timing isolation where the hardware offers the means to
achieve it. Meszlényi et al.\ \cite{timingiso2025} develop timing isolation for
mixed-criticality communication, and Hoornaert et al.\ \cite{mcti2024} present task-based
isolation for mixed-criticality workloads. Burns and Davis's review \cite{mcreview2026}, now in
its fourteenth edition, surveys the field as a whole.

The common prerequisite in this literature is hardware support: memory-bandwidth regulation,
cache partitioning, interrupt steering, or a hypervisor. None of these exists on an ESP32. The
platform offers core affinity and the scheduler's priority mechanism, and nothing else. The
mixed-criticality literature is therefore highly relevant as a statement of what adequate
isolation requires, and inapplicable as a source of techniques for this device class. That
mismatch — a well-developed theory whose preconditions the hardware does not satisfy — is part
of why the standard advice for the ESP32 remains unverified rather than merely weak.

\section{Radio-layer attacks on this device class}

The attack that threatens a wireless alarm is not IP flooding. It is deauthentication: 802.11
management frames handled in the MAC layer and the radio driver, below the IP stack, which the
device must process regardless of what any application is doing.

Kremser and Graffi \cite{deauther2026} evaluate the ESP32-S3 for Wi-Fi penetration testing and
demonstrate that the hardware is itself a capable attack platform, which establishes that the
capability required is cheap and widely available. Cayre et al.\ \cite{espwn2023} examine the
ESP32's own wireless stack offensively, finding vulnerabilities at the system-on-chip level.
Mudraje et al.\ \cite{esp32c3re2025} reverse-engineer the ESP32-C3's Wi-Fi drivers, work
undertaken precisely because the driver is closed source — a fact with direct consequences for
this project, since an unanalysable driver is why one interference term in
Chapter~\ref{ch:impl}'s response-time decomposition cannot be settled analytically.

Carbajal and Jodeiri Akbarfam \cite{carbajal2026} provide a software-defined testbed for
quantifying deauthentication resilience. Their work is methodologically the nearest to the
hardware protocol specified in Appendix~\ref{app:hardware}, and it reinforces a constraint
observed here: experiments of this kind require an isolated radio environment, which is an
institutional rather than a technical prerequisite.

What this literature establishes is that the radio-layer attack path is real, cheap to mount
and distinct from the IP path. What it does not do is measure the \emph{safety} consequence on
a device whose safety task is pinned to a separate core. That is the question this project set
out to answer and, as Chapter~\ref{ch:limitations} states plainly, could not answer by
measurement.

\section{Finding defects in embedded firmware}

Yun et al.\ \cite{fuzzsurvey2022} survey fuzzing for embedded systems and identify the
recurring obstacle: embedded firmware is difficult to execute outside its target hardware, so
fuzzing campaigns must either run on the device, at low throughput, or rehost the firmware,
with fidelity loss.

Two recent systems address rehosting directly. Bley et al.\ \cite{pemu2025} present
protocol-aware firmware rehosting for fuzzing embedded network stacks; Tharindu et al.\
\cite{sysfuss2026} combine system-level firmware fuzzing with selective symbolic execution.
Both are substantially more capable than the approach taken here, and the comparison is worth
stating honestly: the fuzzing in this project compiles the parser \emph{for the host} and
exercises it under sanitizers. That is appropriate for the question asked — whether the command
parser mishandles malformed input — and it is explicitly \emph{not} firmware rehosting. It
says nothing about faults depending on the Xtensa toolchain's code generation, on the real
memory map, or on the radio driver's own parsing. Chapter~\ref{ch:limitations} records this as a
limitation rather than leaving it implicit.

The more important point for this work, however, is one the fuzzing literature does not
address. All of these techniques test a \emph{running system}. None of them tests the
\emph{assumptions a system was built on}. A build condition that encodes an unsound security
inference will compile, run, and pass every dynamic test, because dynamic testing exercises the
build that was produced rather than questioning whether that build should have been permitted.
The unbounded drain reported in Section~\ref{sec:impl-drain} makes the same point from the
timing side: it survived 3,000,000 fuzzing iterations because every one of those iterations
exercised a path that terminated.

\section{Security and verification on the ESP32 application stack}

Azhar et al.\ \cite{mqtttls2026} analyse the performance and security of MQTT over TLS on
ESP32-based systems, establishing the cost of the transport protection this project assumes as
a baseline. Ieropoulos et al.\ \cite{ondevice2026} place intrusion detection on constrained
devices, which is a different architectural answer to a related problem.

On the MQTT protocol itself, Zeghida et al.\ \cite{zeghida2024} give an overview of security
issues, and Alsabbagh et al.\ \cite{pirates2024} demonstrate the attack class that this
project's configuration defect exposes: a rogue client with valid broker credentials raiding
industrial IoT systems. This citation matters for the integrity of the present work's claims.
\textbf{The attack class is published and the underlying principle is textbook}, so this project
claims no new vulnerability. Its contribution is narrower: that a \emph{compile-time guard} can
encode the unsound inference, emit no diagnostic, and thereby cost a safety property.

Two 2026 papers approach the same problem from the design side. Amirkhanova et al.\
\cite{amirkhanova2026} build an end-to-end encrypted data pipeline for industrial IoT, and Yu
et al.\ \cite{mqttl2026} propose MQTT-L, a lightweight end-to-end secure protocol using
renewable hash chains. Both confirm, from the direction of protocol design, the principle this
project confirms from the direction of failure analysis: end-to-end authentication is required
because the broker is a relay and transport security does not survive it.

On formal methods, Krichen's survey \cite{krichen2023} covers verification and validation for
the Internet of Things, including model checking, abstract interpretation and deductive
methods. Lamport \cite{lamport2002} is the reference for the specification language used here.
Neither addresses the failure mode reported in Chapter~\ref{ch:formal} — a specification that
assumes the property it claims to establish — which is a known hazard in the formal-methods
community but, as far as this review found, is rarely disclosed in applied papers.

\section{Research gap}

Five gaps follow from the review above, and they map onto the objectives in
Section~\ref{sec:intro-scope}.

\begin{enumerate}[leftmargin=1.8em,itemsep=5pt]
  \item \textbf{Core pinning has never been evaluated as an isolation defence on this
        hardware.} It is measured as an experimental convenience in \cite{behnke2023} and
        recommended universally in practice, but not tested as a mitigation.
  \item \textbf{Radio-layer interference has not been related to safety outcomes.} The
        degradation literature applies IP load; the attack literature demonstrates radio-layer
        capability; nobody connects the two to a safety function's deadline.
  \item \textbf{The isolation argument has not been treated formally, and no work models an
        adversary} rather than volumetric stress.
  \item \textbf{Mixed-criticality isolation techniques presuppose hardware the ESP32 lacks},
        so the gap cannot be closed by importing them.
  \item \textbf{Existing defect-finding techniques test running systems, not design and
        configuration assumptions.} This is the gap the project's central claim addresses, and
        it emerged from the work rather than preceding it.
\end{enumerate}

\section{How this work differs}

\begin{longtable}{@{}p{0.27\textwidth}p{0.33\textwidth}p{0.33\textwidth}@{}}
\toprule
\textbf{Aspect} & \textbf{Prior work} & \textbf{This work} \\
\midrule
\endhead
Object of study & the phenomenon of degradation \cite{behnke2020} & the \emph{mitigation} that
developers are told to use \\
Load applied & IP layer: UDP, MODBUS/TCP & command interface, and radio layer in the
specified-but-unexecuted protocol \\
Adversary & volumetric stress & an explicit adversary model, scenarios A--E$'$ \\
Isolation argument & informal & TLA+ specification, exhaustively checked, with trace
conformance to the implementation \\
Isolation technique & requires bandwidth regulation, cache partitioning or a hypervisor
\cite{timingiso2025,mcti2024} & core affinity and bounded queues, which is all the hardware
offers \\
Defect finding & fuzzing and rehosting of running firmware
\cite{pemu2025,sysfuss2026} & falsification of build-time, design-time and specification
assumptions \\
MQTT security & protocol design \cite{amirkhanova2026,mqttl2026} & where an existing
configuration guard encodes an unsound inference \\
Timing claim & measured on hardware & \textbf{not measured}; protocol published instead \\
\bottomrule
\end{longtable}

The last row is the one to read carefully. On the question that motivated the project, prior
work has measurements and this work does not. What this work contributes on that question is a
bound, a decomposition that identifies the single term requiring hardware, and a protocol
capable of falsifying the claim. Chapter~\ref{ch:limitations} does not soften this.
